\documentclass{ltjsarticle}
\usepackage{amsmath,amssymb}
\usepackage{amsthm}
\usepackage{amsfonts}
\usepackage{ascmac}
\usepackage{enumerate}
\begin{document}
\theoremstyle{plain}
\newtheorem{prob}{問題}
\newtheorem{lem}{補題}
\newcommand{\Rinf}{\mathbb{R}^{\infty}}
\newcommand{\topoinf}{\mathfrak{O_\infty}}
\newcommand{\mR}{\mathbb{R}}
\newcommand{\mN}{\mathbb{N}}
\newcommand{\fb}{\mathfrak{B}}
\newcommand{\fbb}{\mathfrak{B}'}
\newcommand{\mP}{\mathcal{P}}
%\newcommand{\Im}{\operatorname{Im}}
\newcommand{\Ker}{\operatorname{Ker}}
\newcommand{\rank}{\operatorname{rank}}
\title{代数学 第四回レポート課題}
\author{中野雄太}
\maketitle
\begin{prob}
    完備距離空間$X$と、$X$上の複素数値連続関数の空間$C(X)$を考える。 \\
    $\mathbb{F} \subset C(X)$が同程度連続であるとは、とは
    $\forall x\in X,\forall \epsilon>0,\exists U_x,\forall y\in U_x,\forall$
    この時、同程度連続かつ一様有界であることと、$F$がコンパクトであることは同値である。
\end{prob}